\subsection{有限生成阿贝尔群}

在 \(\mathbf{A} = \mathbf{Z}\) 的情形，第 4 节的结果可以表述为：

定理 3. --- 每个有限生成阿贝尔群 G 都是其挠子群 F（G 中有限阶元素构成的子群）与一个有限秩 p 的自由阿贝尔群（同构于 \(Z^{p}\)）的直和。群 F 是有限多个阶为 \(n_{1}\), \(n_{2}\), \ldots, \(n_{q}\) 的循环群的直和，其中诸 \(n_{i}\) 是 \textgreater1 的整数，且每个整除其前一个；此外，整数 p、q 和 \(n_{i}\)（\(1 \leq i \leq q\)）由 G 唯一确定。

注记. --- 虽然 F 所分解成的那些循环群的阶 \(n_{1}, \ldots, n_{q}\) 由定理 3 中的整除条件完全确定，但这些群本身却并非如此：例如，在 \(\mathbf{Z}/(p)\) 与自身的乘积 G 中（p 为素数），子群恰好是 \(F_{p}\)-向量子空间，而 G 以 \(p(p+1)\) 种不同的方式成为两个一维子空间的直和。

推论 1. --- 在有限阿贝尔群 G 中，存在一个元素，其阶为 G 中所有元素的阶的最小公倍数；这个阶 n 是 G 的第一个不变因子。

推论 2. --- 任何阶不被任何整数 \textgreater1 的平方整除的有限阿贝尔群 G 都是循环群。

让我们沿用定理 3 的记号。于是 \(p = 0\)（因为 \(G\) 有限），且 \(q \leqslant 1\)，否则 \(G\) 的阶会被 \(n_q^2\) 整除。因此 \(G\) 是循环群。

推论3. --- 设 L 和 M 是两个秩为 n 的自由 Z-模，设 \((e_{i})\) 是 L 的一组基，\((f_{i})\) 是 M 的一组基 \((1 \leqslant i \leqslant n)\)，设 u 是从 L 到 M 的一个同态，并设 U 是 u 关于基 \((e_{i})\) 和 \((f_{i})\) 的矩阵。则 Coker \(u = \mathrm{M}/u(\mathrm{L})\) 是有限的当且仅当 \(\det(U) \neq 0\)，且此时 \(\operatorname{Card}(\operatorname{Coker}(u)) = |\det(U)|\)。

如有必要，通过改变 L 和 M 中的基，我们可以假设 U 具有 VII 第 22 页命题 5 的推论 1 中所描述的形式（其中诸 \(\alpha_{i}\) 此时是整数）；该推论随之成为显然，因为 Z-模的直和 \(Z/\alpha_{i}Z\)（\(1 \leqslant i \leqslant n\)）的阶在某个 \(\alpha_{i}\) 为零时是无限的，否则等于 \(|\alpha_{1}\alpha_{2} \ldots \alpha_{n}|\)。

